\section{Pregunta 1. Claves y cifrado de ficheros}
\label{sec:pregunta1}

\begin{enunciado}{(5 puntos)}
Encuentre la mejor solución para dada una clave de usuario de 128 bits generar 10 claves que permitan cifrar 10 ficheros, cada uno con una clave diferente. Implemente la solución para cifrar y descifrar dichos ficheros (su contenido no es relevante).

Pista: Openssl y scripts le ayudarán en esta tarea.
\end{enunciado}

\subsection{Análisis}

El objetivo es obtener diez claves de cifrado diferentes a partir de una clave de usuario de 128 bits, para proteger diez ficheros y recuperarlos posteriormente. Se mantiene como secreto inicial una clave maestra binaria aleatoria $K_M$, de 16 bytes; no se utiliza el contenido de los ficheros para generar las claves.

La derivación de claves se realizará mediante PBKDF2 (Password-Based Key Derivation Function~2) con HMAC (Hash-based Message Authentication Code) y SHA-256 (Secure Hash Algorithm, con salida de 256 bits).
Con la clave simétrica maestra adecuadamente generada, conocer una clave derivada no permite recuperar de forma computacionalmente viable la maestra ni calcular el resto de claves implicadas o generadas.

Esta función KDF (Key Derivation Function) permitirá obtener, para cada fichero $i$, una clave derivada intermedia $\mathrm{DK}_i$ (Derived Key) de longitud $L=48$ bytes a partir de la clave simétrica maestra $K_M$ de 16 bytes, una sal pública aleatoria $S_i$ de 16 bytes (128 bits) distinta para cada fichero y un número de iteraciones $c$ que determina el coste del cálculo. La longitud de la sal reduce la probabilidad de repetición y no depende del tamaño de la maestra, del resumen ni de la salida.

\begin{samepage}
\[
  \mathrm{DK}_i=\operatorname{PBKDF2\text{-}HMAC\text{-}SHA\text{-}256}
  (K_M,S_i,c,L).
\]
El procedimiento se fundamenta en el desarrollo del tema 1.4 y sus diapositivas \ref{ref:tema-kdf}, \ref{ref:slides-kdf}, y se implementará con OpenSSL \ref{ref:openssl-pbkdf2} en un sistema macOS.
\end{samepage}
Se fija $c=600\,000$ iteraciones como coste de derivación. La salida de 48 bytes se dividirá en dos segmentos sin solapamiento: los primeros 16 bytes (128 bits) serán la clave de cifrado $K_i^{\mathrm{cif}}$ y los 32 bytes (256 bits) restantes, la clave auxiliar de autenticación $K_i^{\mathrm{mac}}$. Las diez claves exigidas son las de cifrado; las auxiliares permiten añadir Encrypt-then-MAC (cifrar y después autenticar).

\begin{figure}[H]
  \centering
  \resizebox{\linewidth}{!}{% Original vector diagram: topic 1.4, slide 26 and RFC 8018, section 5.2.
% One 48-byte derivation per file; output slices never overlap.
\begingroup
\setstretch{1}
\hyphenpenalty=10000
\newcommand{\keyblock}[3]{%
  {\fontsize{13}{15}\selectfont\bfseries\boldmath #1}\\[4pt]
  {\fontsize{9}{11}\selectfont\mdseries #2}\\[3pt]
  {\fontsize{8}{10}\selectfont\mdseries (#3)}%
}
\begin{tikzpicture}[x=1cm,y=1cm,
  keynode/.style={proceso,text width=37mm,minimum height=16mm,inner sep=6pt},
  secretkeynode/.style={keynode,fill=white,funcionsecreta},
  keylength/.style={etiqueta,font=\sffamily\fontsize{8}{10}\selectfont,
    inner sep=2pt,fill=white}
]
  \node[secretkeynode] (master) at (0,1.6)
    {\keyblock{$K_M$}{Master key}{shared secret input}};
  \node[keynode,fill=white,funcionpublica] (salt) at (0,-1.6)
    {\keyblock{$S_i$}{Per-file salt}{random public parameter}};
  \node[keynode,text width=43mm,funcioncripto] (pbkdf) at (5.5,0)
    {\keyblock{PBKDF2}{Password-Based Key\\Derivation Function 2}{HMAC-SHA-256;\\600\,000 iterations}};
  \node[secretkeynode] (derived) at (11.4,0)
    {\keyblock{$\mathrm{DK}_i$}{Derived key material}{secret output}};

  \draw[flujo] (master.east) -- node[keylength,above] {16 bytes (128 bits)}
    (5.5,1.6) -- (pbkdf.north);
  \draw[flujo] (salt.east) -- node[keylength,below] {16 bytes (128 bits)}
    (5.5,-1.6) -- (pbkdf.south);
  \draw[flujo] (pbkdf.east) -- node[keylength,above] {48 bytes\\(384 bits)} (derived.west);

  \node[keynode,funcionformato] (split) at (11.4,-2.7)
    {\keyblock{Split}{Fixed byte partition}{non-overlapping slices}};
  \draw[flujo] (derived.south) -- node[keylength,right] {48 bytes (384 bits)} (split.north);

  \node[secretkeynode,text width=43mm] (enc) at (5.5,-5.3)
    {\keyblock{$K_i^{\mathrm{cif}}$}{Encryption key}{bytes 1--16; AES-128}};
  \node[secretkeynode,text width=43mm] (mac) at (11.4,-5.3)
    {\keyblock{$K_i^{\mathrm{mac}}$}{Authentication key}{bytes 17--48; HMAC-SHA-256}};
  \draw[flujo] (split.west) -- node[keylength,above] {16 bytes (128 bits)}
    (5.5,-2.7) -- (enc.north);
  \draw[flujo] (split.south) -- node[keylength,right] {32 bytes (256 bits)} (mac.north);
  \leyendaFunciones{(5.7,-6.7)}
\end{tikzpicture}
\endgroup
}
  \caption{Generación de claves simétricas de cifrado y autenticación mediante PBKDF2 a partir de una clave simétrica maestra única y una sal por fichero. Elaboración propia a partir del tema 1.4 y de la especificación de PBKDF2 (\ref{ref:tema-kdf}, \ref{ref:slides-kdf}, \ref{ref:rfc-pbkdf2}).}
  \label{fig:p1-derivacion}
\end{figure}

PBKDF2 utiliza $K_M$ como clave de HMAC, procesa la sal y aplica el cálculo de HMAC repetidamente, combinando los resultados mediante una operación XOR por bloque de salida en función del número de iteraciones escogido (aumentan el coste del cálculo).
Las sales distintas permiten obtener claves diferentes para cada fichero. La partición asigna bytes diferentes al cifrado y a la autenticación: se realiza una sola derivación de 48 bytes por fichero, no dos derivaciones con la misma sal y longitudes diferentes \ref{ref:rfc-pbkdf2}.

Se conservarán la sal pública por fichero, el identificador, las iteraciones y los algoritmos para reconstruir ambas claves con $K_M$ al descifrar. Los 48 bytes (384 bits) de salida no aumentan la entropía de la clave maestra, que sigue siendo de 128 bits.

El cifrado empleará AES (Advanced Encryption Standard) con claves simétricas de 128 bits (16 bytes) en modo CBC (Cipher Block Chaining), un vector de inicialización aleatorio de 16 bytes por operación y relleno PKCS\#7 (Public-Key Cryptography Standards, number~7), conforme al tema 1.2 \ref{ref:tema-aes}.

Como CBC no detecta por sí mismo modificaciones maliciosas, se añadirá un MAC (Message Authentication Code) mediante HMAC-SHA-256, siguiendo la composición Encrypt-then-MAC estudiada en el tema 1.4 \ref{ref:tema-kdf}.

La clave $K_i^{\mathrm{mac}}$, distinta de la utilizada por AES para cifrar, permitirá autenticar el criptograma, el vector de inicialización y los metadatos de recuperación.

\begin{figure}[H]
  \centering
  \resizebox{\linewidth}{!}{% Original vector diagram: topics 1.2 and 1.4, Encrypt-then-MAC.
% Variable byte lengths: n_i plaintext, m_i padded plaintext/ciphertext,
% h_i serialized header and d_i unambiguous header/ciphertext encoding.
\begingroup
\setstretch{1}
\hyphenpenalty=10000
\newcommand{\cipherblock}[3]{%
  {\fontsize{13}{15}\selectfont\bfseries\boldmath #1}\\[4pt]
  {\fontsize{9}{11}\selectfont\mdseries #2}\\[3pt]
  {\fontsize{8}{10}\selectfont\mdseries (#3)}%
}
\begin{tikzpicture}[x=1cm,y=1cm,
  ciphernode/.style={proceso,text width=37mm,minimum height=16mm,inner sep=6pt},
  publicnode/.style={ciphernode,fill=white,funcionpublica},
  secretnode/.style={ciphernode,fill=white,funcionsecreta},
  datalength/.style={etiqueta,font=\sffamily\fontsize{8}{10}\selectfont,
    inner sep=2pt,fill=white}
]
  \node[secretnode] (plain) at (0,0)
    {\cipherblock{$M_i$}{Plaintext file}{variable length}};
  \node[ciphernode,funcionformato] (padding) at (5.8,0)
    {\cipherblock{PKCS\#7}{Block padding}{add 1--16 bytes;\\$m_i=16(\lfloor n_i/16\rfloor+1)$}};
  \node[ciphernode,text width=40mm,funcioncripto] (aes) at (11.6,0)
    {\cipherblock{AES-128-CBC}{Advanced Encryption Standard\\Cipher Block Chaining}{key $K_i^{\mathrm{cif}}$;\\random vector $V_i$ (16 bytes)}};
  \draw[flujo] (plain.east) -- node[datalength,above] {$n_i$ bytes\\($8n_i$ bits)} (padding.west);
  \draw[flujo] (padding.east) -- node[datalength,above] {$m_i$ bytes\\($8m_i$ bits)} (aes.west);

  \node[publicnode,text width=40mm] (header) at (5.8,-3.3)
    {\cipherblock{$H_i$}{Public header}{version, $S_i$, $c$, file index,\\algorithms and $V_i$}};
  \node[publicnode] (cipher) at (11.6,-3.3)
    {\cipherblock{$C_i$}{Ciphertext}{multiple of 16 bytes}};
  \draw[flujo] (aes.south) -- node[datalength,right] {$m_i$ bytes\\($8m_i$ bits)} (cipher.north);

  \node[ciphernode,text width=40mm,funcionformato] (encode) at (5.8,-6.6)
    {\cipherblock{$D_i$}{Unambiguous encoding}{$D_i=\operatorname{cod}(H_i,C_i)$}};
  \node[ciphernode,text width=40mm,funcioncripto] (hmac) at (11.6,-6.6)
    {\cipherblock{HMAC-SHA-256}{Hash-based Message\\Authentication Code}{key $K_i^{\mathrm{mac}}$; input $D_i$}};
  \draw[flujo] (header.south) -- node[datalength,left] {$h_i$ bytes\\($8h_i$ bits)} (encode.north);
  \draw[flujo] (cipher.south) -- (11.6,-4.95)
    -- node[datalength,above] {$m_i$ bytes ($8m_i$ bits)} (6.65,-4.95)
    -- ([xshift=8.5mm]encode.north);
  \draw[flujo] (encode.east) -- node[datalength,above] {$d_i$ bytes\\($8d_i$ bits)} (hmac.west);

  \node[publicnode] (tag) at (11.6,-9.9)
    {\cipherblock{$T_i$}{Authentication tag}{32-byte HMAC output}};
  \node[publicnode,text width=40mm] (package) at (5.8,-9.9)
    {\cipherblock{$D_i\,\Vert\,T_i$}{Protected file}{encoded data and tag}};
  \draw[flujo] (hmac.south) -- node[datalength,right] {32 bytes (256 bits)} (tag.north);
  \draw[flujo] (encode.south) -- node[datalength,left] {$d_i$ bytes\\($8d_i$ bits)} (package.north);
  \draw[flujo] (tag.west) -- node[datalength,above] {32 bytes\\(256 bits)} (package.east);
  \leyendaFunciones{(5.8,-11.4)}
\end{tikzpicture}
\endgroup
}
  \caption{Cifrado y autenticación de un fichero con el método Encrypt-then-MAC. El vector $V_i$ es aleatorio y la cabecera se autentica junto al criptograma. Elaboración propia a partir de los temas 1.2 y 1.4 (\ref{ref:tema-aes}, \ref{ref:tema-kdf}).}
  \label{fig:p1-cifrado}
\end{figure}

Finalmente, se reconstruirán las claves mediante PBKDF2 con $K_M$, $S_i$ y los parámetros originales, incluidas las iteraciones $c$.
Se recalculará la etiqueta con $K_i^{\mathrm{mac}}$ y se comparará en tiempo constante con la recibida. Si no coincide, se rechazará el fichero al no poder verificar su integridad y autenticidad.
Si coincide, se aplicará el descifrado AES-128-CBC con $K_i^{\mathrm{cif}}$ y el vector de inicialización $V_i$ originales. Tras validar y retirar el relleno PKCS\#7, si precisa, se entregará la información original.

\begin{figure}[H]
  \centering
  \resizebox{\linewidth}{!}{% Original vector diagram: verify Encrypt-then-MAC before decryption.
% Control-flow arrows carry decisions; data arrows carry byte/bit lengths.
\begingroup
\setstretch{1}
\hyphenpenalty=10000
\newcommand{\verifyblock}[3]{%
  {\fontsize{13}{15}\selectfont\bfseries\boldmath #1}\\[4pt]
  {\fontsize{9}{11}\selectfont\mdseries #2}\\[3pt]
  {\fontsize{8}{10}\selectfont\mdseries (#3)}%
}
\begin{tikzpicture}[x=1cm,y=1cm,
  verifynode/.style={proceso,text width=37mm,minimum height=16mm,inner sep=6pt},
  datalength/.style={etiqueta,font=\sffamily\fontsize{8}{10}\selectfont,
    inner sep=2pt,fill=white}
]
  \node[verifynode,fill=white,funcionpublica,text width=40mm] (package) at (0,0)
    {\verifyblock{$D_i\,\Vert\,T_i$}{Protected file}{encoded data and tag}};
  \node[verifynode,funcioncripto,text width=40mm] (check) at (6.4,0)
    {\verifyblock{HMAC-SHA-256}{Hash-based Message\\Authentication Code}{key $K_i^{\mathrm{mac}}$; input $D_i$}};
  \node[verifynode,funcioncripto,text width=40mm] (valid) at (12.8,0)
    {\verifyblock{$T'_i\stackrel{?}{=}T_i$}{Tag verification}{constant-time comparison}};
  \draw[flujo] (package.east) -- node[datalength,above] {$d_i$ bytes\\($8d_i$ bits)} (check.west);
  \draw[flujo] (check.east) -- node[datalength,above] {32 bytes\\(256 bits)} (valid.west);
  \draw[flujo] (package.north) -- (0,2.6)
    -- node[datalength,above] {$T_i$: 32 bytes (256 bits)} (12.8,2.6)
    -- (valid.north);

  \node[verifynode,fill=white] (reject) at (12.8,-4.2)
    {\verifyblock{Reject}{Stop processing}{no plaintext output}};
  \node[verifynode,funcioncripto,text width=40mm] (decrypt) at (6.4,-4.2)
    {\verifyblock{AES-128-CBC}{Advanced Encryption Standard\\Cipher Block Chaining}{decrypt $C_i$; key $K_i^{\mathrm{cif}}$;\\vector $V_i$}};
  \node[verifynode,funcionformato,text width=40mm] (unpad) at (0,-4.2)
    {\verifyblock{PKCS\#7}{Padding validation}{check and remove\\1--16 padding bytes}};
  \draw[flujo] (valid.south) -- node[datalength,right] {No} (reject.north);
  \draw[flujo] (valid.south west) |- (9.0,-2.4)
    -- node[datalength,above] {Yes} (6.4,-2.4) -- (decrypt.north);
  \draw[flujo] (decrypt.west) -- node[datalength,above] {$m_i$ bytes\\($8m_i$ bits)} (unpad.east);

  \node[verifynode,fill=white,funcionsecreta,text width=40mm] (plain) at (0,-7.8)
    {\verifyblock{$M_i$}{Recovered plaintext}{released after validation}};
  \draw[flujo] (unpad.south) -- node[datalength,left] {Valid padding\\$n_i$ bytes ($8n_i$ bits)} (plain.north);
  \draw[flujo] (unpad.south east) -- (1.6,-6.1)
    -- node[datalength,below] {Invalid padding} (12.8,-6.1)
    -- (reject.south);

  \leyendaFunciones{(6.4,-9.2)}
\end{tikzpicture}
\endgroup
}
  \caption{Verificación de la autenticidad del fichero y descifrado del mismo. Las claves se reconstruyen mediante PBKDF2 con la sal y el número de iteraciones conservados. Elaboración propia a partir de los temas 1.2 y 1.4 (\ref{ref:tema-aes}, \ref{ref:tema-kdf}).}
  \label{fig:p1-descifrado}
\end{figure}

La solución separa las claves por fichero y por finalidad, manteniendo un único secreto principal.
Sin embargo, derivar más claves no incrementa la entropía de la clave maestra: su protección sigue siendo esencial, ya que su compromiso afectaría a la protección de todos los ficheros.


\subsection{Implementación y validación}
La solución se implementa mediante tres scripts compatibles con \texttt{/bin/sh} (\texttt{bash}).
Los datos se organizan en el directorio \texttt{crypto-workspace} que se crea donde se ejecuten los scripts, separando los secretos, las sales públicas y los ficheros originales, cifrados y recuperados. Se empleó OpenSSL~4.1.0, instalado mediante Homebrew en macOS:\\

\lstinputlisting[firstline=1,lastline=4,caption={Comprobación del entorno de ejecución.},label={lst:p1-entorno-real}]{recursos/salidas/pregunta-1/entorno.txt}

\textbf{Generación de la clave simétrica maestra y las sales:}
El script \texttt{gen-keys.sh} crea el espacio de trabajo con directorios separados y utiliza \texttt{openssl rand} para generar una clave maestra y diez sales de 16 bytes. 
\texttt{umask 077} restringe el acceso a los archivos nuevos al usuario propietario, evitando que otros usuarios del sistema puedan acceder a ellos y modificarlos.
La función \texttt{generate} llama al comprobador antes de confirmar que la inicialización ha finalizado correctamente.
Si el directorio ya existe o se produce un fallo, su eliminación y el nuevo intento requieren escribir \texttt{DELETE}; cualquier otra respuesta conserva los archivos y cancela la ejecución.

\begingroup
\renewcommand{\lstlistingname}{Script}
\setstretch{1}
\lstinputlisting[style=programacompleto,language=bash,
 caption={Inicialización y comprobación del espacio de trabajo: gen-keys.sh.},
 label={lst:p1-gen-completo}]{recursos/listados/pregunta-1/gen-keys.sh}
\endgroup

La ejecución correcta de \texttt{gen-keys.sh} mostró la siguiente salida, que confirma la validación de la maestra y las diez sales:

\begingroup
\renewcommand{\lstlistingname}{Salida y verificación}
\setstretch{1}
\lstinputlisting[style=terminal,basicstyle=\ttfamily\footnotesize,caption={Ejecución de gen-keys.sh con validación satisfactoria.},label={lst:p1-inicializacion-real}]{recursos/salidas/pregunta-1/inicializacion-validada.txt}
\endgroup

A continuación se muestran la clave maestra y las dos primeras sales en hexadecimal, consultadas con \texttt{od} para verificar su contenido visualmente.
Cada par de dígitos representa un byte y cada línea contiene los 16 bytes del archivo.

\begingroup
\setstretch{1}
\lstinputlisting[style=terminal,basicstyle=\ttfamily\footnotesize]{recursos/salidas/pregunta-1/maestra-y-sales-hex.txt}
\endgroup

\textbf{Validación de la generación de claves:} El script \texttt{prepare-inputs.sh} comprueba que existen tanto la clave maestra como las diez sales, y que cada archivo contiene exactamente 16 bytes.
Después compara los 45 pares posibles de sales mediante \texttt{cmp -s}, evitando comparaciones repetidas.
Un archivo ausente, un tamaño incorrecto o dos sales iguales provocan la finalización con error.
El procedimiento no modifica las entradas y solo muestra \texttt{PASS} cuando todas las comprobaciones se superan.
Este script se ejecuta automáticamente desde \texttt{gen-keys.sh} y no es necesario repetirlo durante la derivación de claves.

\begingroup
\renewcommand{\lstlistingname}{Script}
\setstretch{1}
\lstinputlisting[style=programacompleto,language=bash,
 caption={Validación de la clave maestra y las sales: prepare-inputs.sh.},
 label={lst:p1-prepare-completo}]{recursos/listados/pregunta-1/prepare-inputs.sh}
\endgroup

\textbf{Derivación y separación de claves:} El script \texttt{derive-keys.sh} aplica PBKDF2 con HMAC-SHA-256 y 600\,000 iteraciones para cada sal, manteniendo la misma clave maestra.
Los comandos \texttt{od} y \texttt{tr} representan las entradas binarias en hexadecimal; las opciones \texttt{hexpass} y \texttt{hexsalt} permiten que OpenSSL recupere sus bytes originales \ref{ref:openssl-kdf-cli}.
Cada llamada genera 48 bytes: \texttt{dd} guarda los primeros 16 como clave de cifrado y los 32 restantes como clave de autenticación.

Tras separar cada par, se elimina la clave derivada intermedia y \texttt{wc -c} muestra los tamaños reales de las dos claves generadas para cada fichero.
Las claves se guardan en \texttt{secrets/derived-keys} y \texttt{iterations.txt} conserva el coste empleado (número de iteraciones).
La creación del directorio mediante \texttt{mkdir} impide sobrescribir una derivación existente; \texttt{set -eu} detiene el script ante órdenes fallidas o variables no definidas.

\begingroup
\renewcommand{\lstlistingname}{Script}
\setstretch{1}
\lstinputlisting[style=programacompleto,language=bash,
 caption={Derivación y separación de claves de cifrado y autenticación por cada fichero: derive-keys.sh.},
 label={lst:p1-derive-completo}]{recursos/listados/pregunta-1/derive-keys.sh}
\endgroup

La ejecución de \texttt{derive-keys.sh} produjo la siguiente salida.
Para cada fichero, \texttt{wc -c} confirmó 16 bytes de cifrado y 32 bytes de autenticación, con un subtotal de 48 bytes generados por fichero:

\begingroup
\renewcommand{\lstlistingname}{Salida y verificación}
\setstretch{1}
\lstinputlisting[style=terminal,basicstyle=\ttfamily\footnotesize,caption={Ejecución de derive-keys.sh con validación de tamaños.},label={lst:p1-derivacion-real}]{recursos/salidas/pregunta-1/derivacion.txt}
\endgroup

Como ejemplo, se muestran en hexadecimal las dos claves generadas para el fichero 01.
Cada par de dígitos representa un byte: la clave de cifrado ocupa una línea de 16 bytes y la de autenticación, dos líneas de 16 bytes (32 bytes totales).
La representación permite consultar los valores binarios sin modificar los archivos.

\begingroup
\setstretch{1}
\lstinputlisting[style=terminal,basicstyle=\ttfamily\footnotesize]{recursos/salidas/pregunta-1/claves-derivadas-hex.txt}
\endgroup

El procedimiento completo obtiene, a partir de una única clave simétrica maestra, diez pares de claves destinadas al cifrado y a la autenticación de cada fichero por separado.
La validación confirma los tamaños previstos y la ausencia de sales repetidas, mientras que la conservación de las sales y del número de iteraciones permite reproducir de nuevo la derivación con la misma clave maestra.

